\subsubsection{Miller-Rabin primality}
\begin{lstlisting}[style=mystyle, language=C++]
// TC: O(log^3 n) but with tiny constant
// usa: primalidad de numeros hasta 2^64
#include <bits/stdc++.h>
using namespace std;

using u64 = unsigned long long;
using u128 = __uint128_t;

static u64 modpow_mr(u64 a, u64 e, u64 mod){
	u64 r = 1;
	a %= mod;
	while(e){
		if(e & 1) r = (u128)r * a % mod;
		a = (u128)a * a % mod;
		e >>= 1;
	}
	return r;
}

static bool millerRabin(u64 n){
	if(n < 2) return false;
	for(u64 p : {2ULL, 3ULL, 5ULL, 7ULL, 11ULL, 13ULL, 17ULL, 19ULL, 23ULL, 29ULL, 31ULL, 37ULL}){
		if(n % p == 0) return n == p;
	}
	u64 d = n - 1, s = 0;
	while((d & 1) == 0){ d >>= 1; s++; }
	// bases sufficient for deterministic test below 2^64
	for(u64 a : {2ULL, 325ULL, 9375ULL, 28178ULL, 450775ULL, 9780504ULL, 1795265022ULL}){
		if(a % n == 0) continue;
		u64 x = modpow_mr(a, d, n);
		if(x == 1 || x == n - 1) continue;
		bool comp = true;
		for(u64 r=1 ; r<s ; r++){
			x = (u128)x * x % n;
			if(x == n - 1){ comp = false; break; }
		}
		if(comp) return false;
	}
	return true;
}

int main(){
	// test primality of several numbers up to 2^64
	for(u64 n : {2ULL, 17ULL, 561ULL, 1000000007ULL, (u64)1e18}){
		cout << n << ": " << (millerRabin(n) ? "prime" : "composite") << "\n";
	}
	return 0;
}
\end{lstlisting}
